\section*{第 \ref{ch:g9:landforms} 章——地貌}

\begin{solution}{exo:g9:landforms:1}
(a) 风。(b) 水溶解\omterm{def:g3:stones-pebbles-sand:rock}{岩石}（并在\omterm{def:g9:landforms:karst}{洞穴}中\omterm{def:g5:water-shapes-the-land:deposition}{沉积}方解石）。
(c) 海（\omterm{def:g9:landforms:beach}{沿岸漂移}）。(d) 冰。(e) 海。(f) \omterm{def:g1:landscapes-around-us:water}{河流}。
\end{solution}

\begin{solution}{exo:g9:landforms:2}
(a) \omterm{def:g9:landforms:uplift}{隆升}。(b) \omterm{def:g9:landforms:dune}{滑落面}。(c) \omterm{def:g9:landforms:karst}{落水洞}。(d) \omterm{def:g9:landforms:beach}{沿岸漂移}。
(e) \omterm{def:g9:landforms:karst}{喀斯特}\omterm{def:g1:landscapes-around-us:landform}{地貌}。
\end{solution}

\begin{solution}{exo:g9:landforms:3}
(a) 错误：两角指向下风方向，圆钝的鼻子迎风。
(b) 正确。(c) 错误：风把\omterm{def:g3:stones-pebbles-sand:sand}{沙粒}一粒一粒地搬运，沿缓坡向上，
再沿\omterm{def:g9:landforms:dune}{滑落面}向下。(d) 正确。
\end{solution}

\begin{solution}{exo:g9:landforms:4}
(b)、(d)、(e)、(a)、(c)。
\end{solution}

\begin{solution}{exo:g9:landforms:5}
(a) 风从西方吹来：两角指向下风方向。(b) 位于东侧，
两角之间。(c) $60 \div 6 = 10$ 年。(d) 较大的\omterm{def:g9:landforms:dune}{沙丘}含沙更多，
所有\omterm{def:g3:stones-pebbles-sand:sand}{沙子}都须越过丘脊搬运，
才能使\omterm{def:g9:landforms:dune}{沙丘}向前移动。
\end{solution}

\begin{solution}{exo:g9:landforms:6}
(a) $7 \times 1 = \qty{7}{m}$。(b) $2 \div 12 = \frac{1}{6}$ 年。
(c) 在 $\frac{1}{6}$ 年内后退 7 年：$7 \times 6 = 42$
倍快。(d) 海浪先\omterm{def:g5:water-shapes-the-land:erosion}{侵蚀}出一个缺口；上方\omterm{def:g3:stones-pebbles-sand:rock}{岩石}悬于缺口之上，
暂时尚未移动；随后悬空部分整体崩落，
崖缘随之骤然后退。
\end{solution}

\begin{solution}{exo:g9:landforms:7}
(a) 向东。(b) 沿斜向东北方向的\omterm{def:g9:landforms:beach}{海滩}向上冲三次，
每次之后都直直退回向南：形成
向东移动的之字形。(c) \omterm{def:g9:landforms:beach}{沙嘴}正在向东生长，横跨
海湾口。(d) \omterm{def:g9:landforms:beach}{沙嘴}使其后方水域避开海浪；在平静的
水中，细泥可以\omterm{def:g5:water-shapes-the-land:deposition}{沉积}下来。
\end{solution}

\begin{solution}{exo:g9:landforms:8}
(a) 它从空气和土壤中吸收二氧化碳。(b) 弱酸
缓慢溶解石灰岩，但几乎不\omterm{def:g5:water-shapes-the-land:erosion}{侵蚀}花岗岩。(c) 水
只能通过\omterm{def:g3:stones-pebbles-sand:rock}{岩石}的裂隙并沿岩层进入\omterm{def:g3:stones-pebbles-sand:rock}{岩石}；它首先
拓宽这些通道。(d) 地表以下的\omterm{def:g3:stones-pebbles-sand:rock}{岩石}被溶解，
地面缓慢下沉；或\omterm{def:g9:landforms:karst}{洞穴}顶部坍塌。
\end{solution}

\begin{solution}{exo:g9:landforms:9}
(a) 水滴接触洞内空气时，释放出部分二氧化碳，
不能再容纳全部溶解的石灰岩，于是析出
方解石。(b) 石笋由溅落在
地面并向四周铺开的水滴生长而成；没有水滴悬挂在顶端，以形成管状结构。
(c) $\qty{30}{cm} = \qty{300}{mm}$；$300 \div 0.1 = 3000 = 3 \times 10^3$
年。
\end{solution}

\begin{solution}{exo:g9:landforms:10}
(a) $4600 \div 10^7 = 4.6 \times 10^{-4}$~m 每年，即 \qty{0.46}{mm} 每
年。(b) $0.46 - 0.2 = \qty{0.26}{mm}$ 每年。(c) $\qty{1}{m} =
\qty{1000}{mm}$；$1000 \div 0.46 \approx 2200$，约快两
千倍。(d) 这一速率若持续数百万年：$\qty{0.26}{mm} \times
10^7 = 2.6 \times 10^6$~mm，约 \qty{2.6}{km}。
\end{solution}

\begin{solution}{exo:g9:landforms:11}
(a) 形态：新月形\omterm{def:g9:landforms:dune}{沙丘}，\omterm{def:g9:landforms:dune}{滑落面}位于东南侧。
营力与过程：风从西北吹来，\omterm{def:g5:water-shapes-the-land:deposition}{沉积}作用。
物质：松散的沙。(b) 形态：一道\omterm{def:g9:landforms:cliff}{悬崖}俯临一片平坦台面，
低潮时露出。营力与过程：海，\omterm{def:g5:water-shapes-the-land:erosion}{侵蚀}作用；台面为
随\omterm{def:g9:landforms:cliff}{悬崖}后退留下的海蚀平台。物质：白色\omterm{def:g3:stones-pebbles-sand:rock}{岩石}，可能是白垩，
一种较软的\omterm{def:g3:stones-pebbles-sand:rock}{岩石}。(c) 形态：一条溪流消失在\omterm{def:g9:landforms:karst}{洞穴}中，
远处有泉水涌出。营力与过程：水溶解\omterm{def:g3:stones-pebbles-sand:rock}{岩石}，
并在地下流动。物质：石灰岩；该高原是\omterm{def:g9:landforms:karst}{喀斯特}\omterm{def:g1:landscapes-around-us:landform}{地貌}。
\end{solution}

\begin{solution}{exo:g9:landforms:12}
(a) 小\omterm{def:g9:landforms:dune}{沙丘}：$45 \div 6 = 7.5$ 年；大\omterm{def:g9:landforms:dune}{沙丘}：$45 \div 3 =
15$ 年。(b) $45 \div 100 = 0.45$ 年，约五到六个月。
(c) 例如，在\omterm{def:g9:landforms:dune}{沙丘}上种植草本植物：其根系固定\omterm{def:g3:stones-pebbles-sand:sand}{沙子}，
使风无法再将其带走。
\end{solution}

\begin{solution}{pb:g9:landforms:1}
\textbf{1.} $46.0 - 30.4 = \qty{15.6}{m}$。

\textbf{2.} $15.6 \div 20 = \qty{0.78}{m}$ 每年。

\textbf{3.} 2012 年至 2014 年间：$41.0 - 36.2 = \qty{4.8}{m}$，两年内
共后退，即 $4.8 \div 2 = \qty{2.4}{m}$ 每年。

\textbf{4.} $2.4 \div 0.78 \approx 3$ 倍快。可能是\omterm{def:g5:natural-hazards:storm}{风暴}频发的
冬季，期间大片悬空岩体崩落。

\textbf{5.} \qty{78}{cm} 每年；$0.78 \times 100 = \qty{78}{m}$ 每
世纪。

\textbf{6.} 是：\qty{0.78}{m} 每年略小于每年一米。

\textbf{7.} $17 \div 0.78 \approx 22$ 年。

\textbf{8.} \omterm{def:g9:landforms:cliff}{悬崖}以跳跃方式后退：两年的间隔内可能出现大规模崩落，
也可能毫无变化。在 20 年的时间尺度上，跳跃与平静的年份
相互抵消。

\textbf{9.} 崖缘还需后退 $30.4 - 10 = \qty{20.4}{m}$：$20.4 \div 0.78
\approx 26$ 年，即 2050 年。

\textbf{10.} $2 \times 0.78 = \qty{1.56}{m}$ 每年；$30.4 \div 1.56
\approx 19.5$ 年，自 2024 年起算，即 2043 年。

\textbf{11.} \omterm{def:g9:landforms:cliff}{悬崖}边缘随时可能崩塌；站在崖边
或其附近都是危险的。

\textbf{12.} $30.4 \div 0.78 \approx 39$ 年，自 2024 年起算：在
\textbf{2063} 年。
\end{solution}
